\chapter*{Abstract}\label{abstract}
\phantomsection
\pdfbookmark[1]{Abstract}{abstract}
\prevspace
Long ECG recordings contain many beats, making manual review difficult.
This study developed a single-lead electrocardiogram (ECG) system to detect
abnormal beats and classify them into five supported groups. A binary Extra
Trees model first separates normal and abnormal beats. A second Extra Trees
model assigns an abnormal subtype. The six output codes are N, PVC, PAC,
LBBB, RBBB, and AFib. These are study-specific annotation groups; in
particular, AFib includes atrial flutter, and some groups include escape beats.

A convolutional autoencoder was trained on the first stored ECG channel of
the MIT--BIH Normal Sinus Rhythm Database, using 15 subjects for fitting and
three for validation. Its weights were then fixed. The supervised models
used modified Lead II (MLII) data from 46 MIT--BIH Arrhythmia Database records,
representing 45 subject clusters. Each beat was described by 223 waveform,
spectral, and statistical features, 15 RR-interval features, and four
autoencoder residual features. Supervised windows were filtered and
standardised separately.

The final 20\% of accepted beats in each record formed the temporal test.
Earlier beats were used for model selection and fitting, with a 16-beat gap
before each evaluation boundary. Because the test segment had been viewed
during earlier development, this was a retrospective evaluation of known
patients, with a risk of optimistic performance estimates.

Binary accuracy was 98.81\% on 20,969 beats, and conditional subtype accuracy
was 97.63\% on 7,958 supported abnormal beats. Whole-system accuracy was
98.07\% and macro F1 was 0.9524 on 20,043 supported beats. Mean accuracy
across subjects was 98.03\%, with a subject-clustered 95\% confidence interval
of 96.29--99.25\%. However, PAC recall was only 70.63\%. A limited external
check used 5,416 beats from two INCART records belonging to one patient.
Whole-system accuracy was 92.85\%, while binary precision fell to 64.22\%.

The system met the aggregate accuracy and interval-width targets in the
retrospective test. It did not meet the recall target for every class.
Independent patient evaluation and real-sensor testing are still needed.
